\section{LITERATURE REVIEW}

\subsection{Execution-Grounded LLM Algorithm Design}

Large language models have increasingly moved from one-shot solution
generation toward feedback-driven optimization. OPRO treats an LLM as a
black-box optimizer by conditioning new proposals on previously evaluated
solution--score pairs~\cite{yang2024opro}, while Self-Refine demonstrates how
verbal feedback can support iterative revision without additional
training~\cite{madaan2023selfrefine}. For executable algorithm discovery,
FunSearch couples an LLM with deterministic evaluation and an evolutionary
program database~\cite{romera2024funsearch}. Evolution of Heuristics evolves
both heuristic ideas and executable implementations~\cite{liu2024eoh}, and
ReEvo converts performance comparisons into reflective feedback for subsequent
generations~\cite{ye2024reevo}. MCTS-AHD organizes candidates in a search tree
to improve exploration~\cite{zheng2025mctsahd}, whereas AlphaEvolve scales the
same generate--evaluate--retain principle through an evolutionary coding-agent
pipeline with automated evaluators~\cite{novikov2025alphaevolve}. Collectively,
these studies show that LLM proposals become more reliable when validity and
quality are determined by external execution rather than linguistic judgment
alone. Our framework adopts this principle, but its search objects are
structured traffic interventions and its evaluator is a microscopic traffic
simulator.

\subsection{Multi-objective, Compositional, and Multi-agent Search}

Most LLM-based algorithm-design systems search for a single program or
heuristic under a scalar score. MEoH extends this setting by maintaining
non-dominated heuristics and accounting for both objective-space dominance and
program diversity~\cite{yao2025meoh}. E2OC addresses interdependent operator
design by representing composition as a sequential decision process and using
Monte Carlo tree search~\cite{qiu2026e2oc}, while OptiMUS illustrates a
complementary form of execution grounding in which LLM-generated optimization
models are checked and solved by mathematical-programming
solvers~\cite{ahmaditeshnizi2024optimus}. Multi-agent debate further suggests that
role separation and critique can improve proposal diversity and reasoning
quality~\cite{du2024debate}. In our framework, signal, network, and routing
agents therefore generate complementary interventions, but neither the agents
nor the advisory judge determine the final plan. Valid single interventions
are evaluated independently, promising candidates are combined subject to
compatibility and evaluation-budget constraints, and the baseline, singles,
and bundles form a common candidate pool. A constraint-aware selector is used
for the primary experiment, while an offline Pareto audit exposes trade-offs
among traffic metrics, following the general non-dominated-selection principle
used in multi-objective optimization~\cite{deb2002nsgaii}.

\subsection{LLMs for Traffic Control and Simulation-Grounded Optimization}

Before LLM-based control, learning-based traffic-signal research developed
simulator-trained policies such as the max-pressure-inspired
PressLight~\cite{wei2019presslight} and the graph-attention-based
CoLight~\cite{wei2019colight}. Recent LLM approaches instead use language models for
traffic reasoning and decision support. LLM-Assisted Light augments an LLM with
traffic perception and control tools~\cite{wang2024llmassistedlight}; LLMLight
directly selects signal phases from structured
observations~\cite{lai2025llmlight}; and CoLLMLight extends this paradigm to cooperative,
network-wide control~\cite{yuan2026collmlight}. Hybrid methods such as iLLM-TSC
inspect and correct reinforcement-learning actions under abnormal or partially
observed conditions~\cite{pang2024illmtsc}, whereas EvolveSignal searches for
executable signal-control programs using simulator
feedback~\cite{wang2025evolvesignal}. A related line of work uses LLM agents to automate
traffic-simulation workflows: ChatSUMO translates natural-language requests
into SUMO scenarios~\cite{li2024chatsumo}, ChatSUMO-Agent adds closed-loop
planning, tool execution, and refinement~\cite{li2026chatsumoagent}, and the
concurrent TrafficSimAgent preprint studies hierarchical agents for general
traffic-simulation tasks~\cite{du2025trafficsimagent}. In contrast, our focus is
not scenario generation or stepwise signal control, but simulation-grounded
search over heterogeneous network, routing, and signal interventions. SUMO
provides the executable microscopic evaluator~\cite{lopez2018sumo}, while the
cross-layer formulation is motivated by the established coupling between
signal control and route assignment~\cite{meneguzzer1997combined} and between
network design and travelers' route choices~\cite{yang1998networkdesign}.
